\section{Introdução}

Em agentes de programação, a taxa de resolução informa se uma tarefa foi concluída segundo o verificador, mas, sozinha, não informa quanto foi necessário gastar para obtê-la. Repetições, chamadas adicionais e revisão podem alterar simultaneamente a probabilidade de sucesso e o consumo de recursos. Kapoor et al. identificam a ênfase exclusiva em acurácia como uma limitação da avaliação de agentes e formulam a otimização conjunta de acurácia e custo; seus experimentos mostram que considerar custo pode mudar a interpretação de estratégias que parecem superiores quando se observa apenas a taxa de sucesso \cite{Kapoor2024AIAgentsMatter}. Para decisões de engenharia, portanto, a comparação precisa mostrar a qualidade alcançada junto com os recursos consumidos, em vez de tratar qualquer ganho de resolução como vantagem suficiente.

A necessidade de avaliar qualidade e recursos conjuntamente torna o harness parte da condição experimental. Em um estudo controlado com 50 tarefas do Terminal-Bench Pro e os modelos Qwen 3.6 Plus e MiniMax M2.5, as diferenças pareadas de taxa de resolução entre harnesses variaram de 0 a 8 pontos percentuais. Os intervalos de confiança incluíram zero na maioria das comparações, o que não demonstra equivalência entre os harnesses. Ao se considerar o consumo de tokens por tarefa resolvida no mesmo estudo, observaram-se diferenças de até 40 vezes. Esse contraste reforça que a taxa de resolução, isoladamente, omite diferenças substanciais de consumo mesmo em um contexto experimental controlado \cite{Vats2026ScaffoldEffect}. Como a avaliação abrangeu apenas dois modelos e três harnesses, o fator de 40 vezes demonstra a possibilidade de grande variação, mas não estima a magnitude do efeito em outros modelos, sobretudo modelos mais recentes ou mais capazes. Tokens tampouco equivalem a um custo monetário fixo, pois tarifas, cache e categorias de cobrança variam.

Além de afetar custo e desempenho, a infraestrutura pode determinar se uma alteração produzida pelo agente chega sequer a ser testada. O Claw-SWE-Bench evidencia esse problema na interface entre agente e avaliador. Em sua configuração mínima, o modelo precisava escrever um \emph{diff} unificado na resposta final; erros em caminhos, cabeçalhos ou linhas de contexto podiam impedir a aplicação do patch antes da execução dos testes. Na configuração completa, o agente editava os arquivos do repositório, e a camada de integração calculava o \emph{diff} em relação ao estado inicial, removia artefatos alheios à solução e produzia o patch no formato aceito pelo avaliador. Com o mesmo modelo GLM 5.1 e as mesmas 350 tarefas, a taxa de resolução passou de 19,1\% para 73,4\%, enquanto a proporção de patches não aplicáveis caiu de 69,1\% para menos de 1,5\% \cite{Zheng2026ClawSWEBench}. O contraste não atribui esse ganho apenas à extração do patch, pois a configuração completa também modificou outros elementos do adaptador. Ele demonstra, contudo, que a pontuação pode confundir a capacidade de resolver a tarefa com a capacidade de satisfazer o contrato de entrega. A unidade de comparação deve, assim, identificar modelo, harness, ferramentas, orçamento e procedimento de entrega.

Mesmo uma comparação controlada pode ser inválida se as tarefas avaliadas não medirem capacidade de resolver problemas novos. No SWE-bench Verified, um estudo pediu a dois modelos Claude que localizassem arquivos relevantes usando o texto da issue, com contexto deliberadamente limitado, e comparou esse desempenho ao de BeetleBox e SWE-rebench. Os modelos tiveram desempenho três vezes maior no Verified e seis vezes maior na identificação de arquivos editados sem contexto adicional; os autores interpretam a diferença como compatível com exposição prévia às tarefas, mas o experimento não revela diretamente os dados de treinamento nem prova memorização de soluções específicas \cite{Prathifkumar2025ModelMemory}. Para a validade da avaliação, isso sustenta registrar a idade e a origem das tarefas, separar desenvolvimento de avaliação e restringir a interpretação de resultados em benchmarks públicos; não autoriza atribuir automaticamente diferenças entre conjuntos à contaminação.

O custo de avaliar muitas combinações de modelos, harnesses e benchmarks também condiciona o desenho do estudo. Ndzomga examina a redução do conjunto de tarefas por dificuldade histórica e encontra uma faixa intermediária útil para preservar a ordenação de agentes em vários cenários, enquanto a estimativa de pontuações absolutas se mostra mais sensível, especialmente quando muda o scaffold \cite{Ndzomga2026EfficientBenchmarking}. Esse resultado fundamenta uma etapa econômica de triagem, mas não transforma uma amostra escolhida para preservar rankings em estimativa confirmatória da taxa de resolução no benchmark inteiro. O protocolo deste estudo separa essas finalidades e reserva a avaliação de qualidade para tarefas que não tenham sido usadas para ajustar ou selecionar configurações.

A validade de um benchmark envolve ainda a correspondência entre o que ele pontua e o trabalho que se espera de um agente de engenharia de software. Gorinova et al. argumentam que benchmarks de programação podem confundir o efeito do modelo com o do harness, vincular a correção à solução de referência e fornecer pouco sinal sobre a origem das falhas \cite{Gorinova2026MisalignedBenchmarks}. Por se tratar de um artigo de posição, essas críticas não quantificam a prevalência de cada problema nos conjuntos usados aqui. Elas apontam, porém, limites de interpretação que a avaliação pode enfrentar com controles de componentes, registro de falhas e inspeção independente de uma amostra de entregas, mantendo o veredito oficial do benchmark como medida operacional principal.

Resultados fora do domínio de programação reforçam que o efeito de um scaffold pode depender do modelo e da tarefa. No estudo controlado de GAIA, os efeitos variaram entre famílias de modelos e níveis de dificuldade; a expectativa de que o scaffold se tornaria sistematicamente menos relevante com modelos mais capazes não se confirmou no nível mais difícil avaliado \cite{Starace2026ScaffoldGAIA}. Como esse benchmark mede tarefas gerais com ferramentas, e não resolução de issues de software, o resultado não antecipa a direção do efeito em engenharia de software. Ele mantém em aberto, como hipótese exploratória, se a capacidade do modelo modera o benefício do scaffold e justifica medir essa interação em vez de pressupor uma tendência uniforme.

Em conjunto, essas evidências sustentam a tese metodológica deste artigo: a avaliação de agentes de programação deve comparar qualidade sob limites explícitos de custo e tempo, controlar o harness como parte do sistema e tratar a validade das tarefas como condição para interpretar o resultado. A arquitetura em estudo explora uma hipótese ainda não confirmada: distribuir planejamento, execução, escalonamento e, quando viável, subtarefas paralelas entre modelos de custos distintos pode melhorar essa relação. As seções seguintes especificam essa configuração e um protocolo para confrontá-la com agentes de modelo único e controles de repetição ou escalonamento, sem antecipar resultados experimentais.
